\cleardoublepage

\chapter*{Resumen\markboth{Resumen}{Resumen}}

La planificación de cobertura radioeléctrica en redes 5G tácticas ---desplegadas en
escenarios de emergencia, operaciones militares o eventos temporales--- requiere
herramientas que permitan estimar y visualizar la propagación electromagnética sobre el
terreno. Las condiciones de licencia de las soluciones comerciales dependen del producto
y la modalidad contratada, por lo que no se atribuye aquí un coste anual único.

El presente Trabajo de Fin de Grado diseña e implementa un prototipo web para un gemelo
digital de cobertura 5G en la banda FR1 (sub-6~GHz) con una antena, con capacidades descriptivas y de
predicción interactiva inicial. Los escenarios de referencia se generan en MATLAB con
Longley--Rice y declaran \texttt{SRTM\_90m\_global}. Para las nuevas peticiones web, un
motor Java integrado en Spring Boot~3.2 (Java~17) calcula FSPL más difracción de Deygout
sobre un GeoTIFF local cuya procedencia debe documentarse. Una API REST
documentada con OpenAPI conecta ambos flujos con una aplicación Angular~21 que permite
configurar, visualizar y comparar escenarios sobre mapas Leaflet.

La validación cruzada controlada de Java frente a una referencia MATLAB con el mismo
algoritmo FSPL+Deygout y supuestos comunes obtuvo diferencias inferiores a
$10^{-8}$~dB en dos perfiles sintéticos. Confirma la concordancia numérica para esos
casos, pero no es validación de campo. En una ejecución anterior se compararon tres
escenarios MATLAB históricos ---suburbano, rural y urbano--- con recálculos Java
en coordenadas comunes. Los RMSE observados variaron entre 12,242 y
19,830~dB, pero la comparación no aísla los modelos: difieren terreno, balance de potencia
y normalización, y los JSON no se regeneraron tras corregir los metadatos de cuadrícula.
Esa comparación no se pudo repetir en la integración actual sin el GeoTIFF histórico
y no constituye validación con medidas de campo. En el candidato integrado, Maven
registró 82 casos: 79 ejecutados y superados y tres omitidos, sin fallos. El frontend
superó nueve pruebas y su compilación fue correcta. Una prueba funcional local anterior en
navegador con seis JSON aislados verificó la creación y visualización del mapa, el borrado
confirmado y cancelado, JSON inválidos, respuestas GET tardías o fallidas y limpiar una
comparación A=B; no es una suite de navegador en CI ni una medición de latencia.
En la integración reorganizada se comprobó también el recorrido real del navegador
con un GeoTIFF y datos sintéticos aislados. La construcción y ejecución de Docker
quedan pendientes por falta de espacio.

La calibración con observaciones compatibles y la recomendación de emplazamientos se
han probado con datos sintéticos, sin validar su precisión real; la selección voraz no
garantiza un óptimo global. La sincronización con una burbuja física sigue pendiente.
MATLAB continúa siendo necesario para generar y reproducir los
escenarios de referencia, pero no para solicitar una predicción mediante el motor Java.
Los resultados son un proxy de potencia, no una medida física de SS-RSRP, y deben
interpretarse con criterio radioeléctrico.

\vspace{1em}
\textbf{Palabras clave:} gemelo digital, 5G NR, cobertura radioeléctrica, Longley-Rice,
SRTM, Spring Boot, Angular, Leaflet, RSRP, redes tácticas.
